% The data, code and proofs statement (loop0009 item 01). It replaces the
% title-page footnote and every file path and `python -m` command that the
% body used to carry. Those commands, kept here for later items and recorded
% in the paper1/*.md documents that are the sources of record:
%   Section 3 figures and numbers  python -m paper1.section2_figures [--numbers]
%                                  (paper1/relevance.py, paper1/trends.py; offline,
%                                  from the OpenAlex caches of 29 September 2026)
%   Tables (generated)             python -m paper1.latex.make_tables
%   Band values (brute force)      python -m paper1.complex_check
%   Figure 4.1                     python -m paper1.counterexample_figure
%   Validation counts, Section 4.4 paper1/solver_fix.md items 01-05;
%                                  paper1/revised_algorithm.md section 4.4.4
%   Prior reports                  python3 -m paper1.prior_art_sweep
%                                  (paper1/prior_art_counterexample.md)
%   Split re-certification         python -m paper1.solver_fix_split --summary
%                                  (Lean: Search/Split.lean)
%   Dataset                        python -m paper1.dataset (builder);
%                                  python paper1/dataset_check.py --sources (checker)
%   Axiom audit                    lake env lean paper1/axiom_check.lean
\section*{Data, code and proofs}
\label{sec:data}
\addcontentsline{toc}{section}{Data, code and proofs}

Our code, proofs and data are openly available at
\url{https://github.com/Alex-Linhares/MOSP_trough_SAT}; the code and the proofs
are open source under the MIT licence. Everything this paper reports can be
re-run from there, and most of it can be checked
without re-running. The Lean proofs check in one build. Witnesses, minor
certificates and the 34-customer instance have independent checkers, and
refutations up to about 75 customers carry certificates
(Section~\ref{sec:certs}). The largest refutations demand enormous certificates and, hence, can only be
re-run, at up to 2,231 core-hours each
(Section~\ref{sec:practice}). The repository holds:
\begin{itemize}
\item the Lean~4 development, with every theorem the paper names
  (Appendix~\ref{app:lean}) and an axiom check that covers all of them;
\item the source of both our solvers, the open-stacks customer search and the
  graph pathwidth solver, in C and in Python, with our corrected rules as the
  default and the rules of the previous algorithm \citep{ChuStuckey2009} as an
  option;
\item the graph pathwidth solver as a standalone package
  (\repofile{pathwidth_solver/}), with the scripts that produced
  Table~\ref{tab:graphs} and the published values it is compared with;
\item the dataset of Section~\ref{sec:dataset}, the certificates of
  Section~\ref{sec:certs} and the independent checkers of both; the benchmark
  graphs keep the terms of their original distributions;
\item the 34-customer instance of Figure~\ref{fig:wrong}, with the optimum our
  implementation of the previous algorithm reports with the better move off and
  its order, an optimal order, a checker that uses only the Python standard
  library, and the DRAT refutation of 5 stacks with the script that checks it
  (\repofile{paper1/data/false_refutation/});
\item the bibliometric data of Section~\ref{sec:names}: the cached OpenAlex
  responses, the relevance labels and the mechanical rule;
\item a script for every table, figure and number in the paper, each of which
  regenerates it from the recorded data.
\end{itemize}
